\scanpage{170}
de groupe un groupe constant $\pi_S$, alors
\uline{tte composante connexe} $P^0$ se surjecte sur $S$.
Or c'est trivial, car $P$ est \uline{trivial} !

Ceci étant, comme $P^{(0)}\subset\deleted{\unsure{P|\bigcap}}P^{(i)}$ pour tt $i$,
on voit que $\bigcap_i P^{(i)}$ \deleted{\unsure{contient}} un pt
de chaque fibre. \deleted{\unsure{Comme} les inclusions \unsure{de}}
\deleted{\unsure{Il s'ensuit que}}
\deleted{\unsure{fibré des|finie}} \inserted{chaque} $P^{(i)}$ est \unsure{une réunion|un réunion} de
composantes connexes de $P$, il en est de
\unsure{$=$|même} de l'intersection, \deleted{qui} $P'$, qui contient
$P^0$. Elle est stable \unsure{sous|sur} $\pi'$, et on \unsure{conclut|constate}
aussitôt (en regardant sur chaque fibre
géométrique, et utilisant que $P'_{\unsure{\bar s|s}}\ne\varnothing$]
que $P'$ est fibré principal homogène \unsure{sur|sous} $\pi'$.
cqfd.

\uline{Conséquences} Si $\pi$ varie dans la \inserted{\unsure{une|une sous-}} catégorie
des groupes artiniens, \deleted{\unsure{ou si $S$ est}}
stable par produits finis et sous-groupes,
ou si $\pi$ varie dans n'importe quelle
catégorie stable par ces opérations, mais $S$
\unsure{étant|etant} loc noeth --- alors le foncteur
$\pi\leadsto\unsure{\Pi|\pi}^{1}(S,\xi;\pi)$ est strictement proreprésentable\unsure{$^{*}$|/}

% margin note written sideways
\margin{\unsure{$^{*}$} par un progroupe $\widetilde{\pi}_1(S,\xi)$ [qui ne peut plus être interprété comme un simple groupe topologique \unsure{---}]}
% archive sheet number: 169
